% Kapitel 8 -- Von der Schwarzen Frau von Hardegg

\section{Von der Schwarzen Frau von Hardegg}

\mylettrine{E}{s} lebte zu jener Zeit in \gls{Hardegg}\ eine Frau von stiller und gottseliger Art, die der Priester des Ortes im Scherz \gls{Regina}\ nannte; und wie es oft geschieht, blieb an ihr haften, was zuerst nur leichthin gesprochen worden war. Das Volk übernahm den Namen, und bald wußte man sie kaum anders zu nennen. Ihr Wesen war schlicht, ihre Rede selten, ihr Wandel ehrbar, und doch lag in ihr eine solche Festigkeit, dass viele in ihr ein Beispiel rechter Frömmigkeit erkannten.

Der Graf \gls{Kuonrat von Hardegg}\ vernahm von ihr, trat in die Kirche, wohl mehr aus Neugier als aus Andacht, und sah sie mit eigenen Augen. Man sagt, sein Blick blieb länger an ihr haften als an Altar, Kerzen und Meßgewand; daraus wuchs in ihm kein frommer Sinn, sondern unreines Begehren. Er begehrte, sie in die Burg zu bringen, doch sie, die den Glanz weltlicher Herren gering achtete und Gott und der Jungfrau Maria nahestand, wies jede heimliche Werbung ab, sooft man ihr davon zu reden suchte.

Solches Gebaren erzürnte den Herrn. Er ließ den jungen \gls{Kuonrat von Steinare}\ vor sich treten und sprach, wie man erzählt, mit kurzer und harter Stimme: \enquote{Reite hin und bringe mir die Frau. Wenn sie zögert, so nimm zwei Knechte mit. Mir ist nicht um ihre Meinung zu tun.} Da hub \gls{Kuonrat von Steinare}\ an, mit jenem verworrenen Eifer zu reden, der ihm eigen war, sooft er von seiner \gls{Freiheit}\ sprach, und man legte ihm Worte in den Mund wie diese: \enquote{Herr, ein Mann mag Euch dienen mit Schild, Lanze, Sporn und müdem Rücken, doch nicht mit dem innersten Willen, denn der ist kein Sattelzeug, das ein anderer umschnallt. Soll ich reiten, so reite ich, soll ich wachen, so wache ich, soll ich mit der Lanze voraus, so will ich es noch eher tun als zehn andere; doch eine Frau, die nicht gehen will, an ihrer eigenen Weigerung packen und schleppen, das ist ein Handel für Wolfszähne und Kettenhunde, nicht für mich. Heute heißt Ihr mich, ihre Schritte zu binden, morgen heißt mich ein anderer Herr meinen Mund zu verleihen, übermorgen will ein dritter noch die Gedanken in meinem Schädel verzinsen wie einen Acker. So schneidet man einem freien Mann Scheibe um Scheibe aus der Haut und nennt es Dienst. Ich aber sage, die \gls{Freiheit}\ ist ein wunderliches Tier: sperrt man sie in den Stall, so schlägt sie aus; läßt man sie los, so rennt sie in Korn und Dornen und weiß selbst nicht, wem sie schadet; und doch bleibt sie \enquote{Freiheit} und nicht Strick. Wenn die Frau selber käme, wollte ich ihr den Weg weisen, wollte ihr Pferd führen, ja wollte noch die Hunde fernhalten. Soll ich sie aber gegen ihren Sinn herbringen, so muß ich erst meinen eigenen Sinn erwürgen; und wer den eigenen Sinn erwürgt, der taugt bald auch zum Kriechen unter fremden Tischen. Nein, Herr, dazu tauge ich nicht. Lieber nennt mich trotzig, faul, närrisch oder einen schlechten \gls{knappe}n, als dass ich mit fremder Gewalt eine fremde Frau an Euren Herd ziehe.} Ob er alles so sprach, vermag ich nicht zu beschwören; wohl aber steht fest, dass er sich weigerte, und daraus wuchs des Grafen Argwohn nur desto mehr. Und dies will ich ihm, bei all seinen Lastern und bei aller Verworrenheit seiner Rede, nicht verkürzen: er tat hier das Rechte, alldieweil er sich weigerte, an einer wehrlosen Frau fremde Gewalt zu üben.

Als Strafe hieß der Graf ihn, mit dem fahrenden \gls{Gruober}\ nach Böhmen zu reiten und ihm bei seinem schändlichen Geschäften beizustehen; zugleich ließ er \gls{Regina}\ gefangen nehmen. Man berichtet, dass \gls{Gruober}\ und seine Gesellen bei der Tat halfen, denn alles ging leise und schnell wie ein nächtlicher Raubzug. Die Frau aber widersetzte sich mit solcher Kraft des Herzens, dass selbst die Knappen, die sie hielten, vom Schrecken ergriffen wurden.

Der Graf, von zorniger Verblendung beherrscht, gebot darauf, \gls{Regina}\ lebendig in die Felswand gegenüber der Burg einzumauern, auf dass sie in Klage, Hunger und Gram zugrunde gehe. Man trieb sie in eine schmale Nische, schloß Steine und Mörtel hinter ihr und ließ ihr Antlitz der Burg zugewandt, so dass sie noch die Zinnen sehen mußte, unter denen ihr Unheil beschlossen worden war. So, mit zurückgeneigtem Haupt und dem Blick auf die Stätte ihres Verderbens, soll sie elend verschieden sein, wie man sich zu erzählen pflegt.

Der Graf aber wollte sich des grausamen Werkes selbst versichern und ritt zum Felsen, um mit eigenen Augen zu schauen, wie die Strafe vollzogen worden sei. Als er dem Abhang nahte, so heißt es, strauchelte sein Schimmel. Der Herr stürzte in den gähnenden Abgrund, und niemand fand seinen bleichen Leichnam; auch das Ross ward nie mehr gesehen. Die Leute in der Gegend deuteten dies bald als Gottesgericht, bald als Werk des Teufels, und manche sagten, die Strafe sei auf den Grafen selbst zurückgefallen.

Seit jener Zeit erzählten Wanderer, bei Vollmond erscheine auf dem Felsen ein Reiter, der dem Absturz entrinnen wolle und doch immer wieder hinabgerissen werde. Er naht dem Rand, sucht Halt, schlägt mit dem Pferde aus und fällt von neuem, als müsse er dieselbe Stunde ohne Ende erleiden; so wie Sisyphos seinen Felsen den Berg hinauf rollt, nur um ihn wieder zu sehen wie er runter rollt. Auch die Gestalt der toten \gls{Regina}, so meinten einige, sei nicht zur Ruhe gekommen. Bei Nacht, so heißt es, gehe sie zum Kirchhof und suche ihr Grab, klagend und rastlos, denn wer lebendig in Stein beschlossen ward, finde nicht die Ruhe des christlichen Todes.

Der Erzähler überliefert dies mit Scheu vor endgültigem Urteil; ob hier Wunderzeichen oder Einbildung der Leute walten, mag der Leser selbst ermessen. Ich schreibe nur nieder, was man mir berichtet hat: dass die Burg von jener Nacht an unter einem dunklen Namen litt und dass man bei Vollmond Klagen und Hufscharren am Felsen vernahm.

So fügt die Kunde von der Schwarzen Frau von Hardegg der Chronik ein finsteres Exempel hinzu. Denn wer Gewalt an den Schwachen übt und seine Macht für ein Recht hält, der lädt Schuld auf sich, die weder Mauer noch Tiefe verschlingen. Darum blieb in Hardegg nicht nur die Erinnerung an \gls{Regina}, sondern auch die Furcht vor jenem Reiter, der bei hellem Mond noch immer seinem Sturz entgegenreitet.
